\documentclass[10pt]{article}

% --- Preamble ---
\usepackage[utf8]{inputenc}
\usepackage[margin=0.5in]{geometry}
\usepackage{graphicx}
\usepackage{amsmath}
\usepackage{amssymb}
\usepackage{booktabs}
\usepackage{hyperref}
\usepackage{xcolor}
\usepackage{caption}

% Setup hyperref
\hypersetup{
    colorlinks=true,
    linkcolor=blue,
    filecolor=magenta,      
    urlcolor=cyan,
    pdftitle={Thalamocortical Expansion Report},
}

\title{Thalamocortical Expansion and Representation Learning in the Visual System}
\author{Sakin Kirti \\ \small Dippopa Lab Rotation Project \\ \small Winter 2026 \\ \small Department of Neurobiology, UCLA}
\date{\today}

\begin{document}

\maketitle

\begin{abstract}
The conversion of sensory inputs into robust cortical representations is a fundamental operation of the mammalian brain. In the early visual system, this process 
is characterized by a dramatic expansion in neuronal density as signals traverse from the Lateral Geniculate Nucleus (LGN) to the Primary Visual Cortex (V1). 
This ``thalamocortical expansion'' has long been hypothesized to facilitate the creation of overcomplete representations, yet the functional advantages of such 
high-dimensional coding remain difficult to quantify at scale. This report investigates the role of expansion in representation learning using a biologically constrained 
feedforward architecture. By simulating populations across a wide range of expansion factors, we demonstrate that a massive cortical over-representation is critical for 
the emergence of sparse, selective features that efficiently separate temporal variations in natural stimuli. Our findings provide support for the hypothesis that the 
thalamocortical expansion serves as a high-dimensional mapping stage necessary for complex metric learning in the visual hierarchy.
\end{abstract}

\section{Introduction}
The architectural layout of the early visual pathway provides a compelling example of specialized biological wiring, characterized by a dramatic transformation of signal 
density as information traverses from the retina to the cortex. Visual signaling begins with a mosaic of Retinal Ganglion Cells (RGCs) that form a functional bottleneck, 
projecting through the optic nerve to the Lateral Geniculate Nucleus (LGN) with a nearly 1:1 ratio. However, the subsequent projection from the LGN to the Primary Visual 
Cortex (V1) undergoes a radical expansion, governed by a power law relationship roughly expressed as $K=N^{3/2}$, where $K$ and $N$ represent the number of V1 and LGN 
neurons respectively. In human systems, this mapping reaches an extreme 1:300 ratio, concentrating a relatively small number of thalamic inputs onto a vast cortical population.

The functional necessity of this massive over-representation is a central theme in computational neuroscience, primarily hypothesized to enable the cortex to implement 
an overcomplete basis set for signal representation. By projecting low-dimensional LGN outputs into a much higher-dimensional space in V1, the brain may perform a 
non-linear mapping that facilitates efficient coding and simplifies the separation of overlapping stimulus features, similar to a kernel expansion in machine learning. 
This structural expansion follows precise topographic and probabilistic wiring rules, where V1 neurons integrate information from localized neighborhoods to diversify 
the feature space while maintaining a retinotopic map. This project aims to quantify the benefits of this expansion by assessing how varying expansion factors directly 
influence the quality of representation learning and temporal separation when processing naturalistic video stimuli.

\section{Methodology}

The trajectory of this project involved an iterative refinement of the model architecture to best capture biological principles while remaining computationally tractable 
for large-scale benchmarking. 

Initially, effort was directed toward implementing the early visual system model proposed by Slapik \& Shouval (2026). While the Slapik model offers insights 
into estimated representations created by each cell type, it relies heavily on a non-feedforward architecture, where each cell type is effectively implemented as a separate 
one-layer feedforward network. During development, it was determined that non-feedforward nature diverged from the primary goal of studying the feedforward expansion bottleneck.  

Consequently, we pivoted to the Ringach (2004) model. The Ringach architecture provides a strictly feedforward framework that more directly isolates the effects of the LGN-to-V1 
expansion. By focusing on a differentiable, feedforward pass, we were able to leverage modern deep learning frameworks like JAX to simulate V1 populations 
in the hypothesized probabilistic, fully learned, and distance-regularized wiring frameworks while maintaining a clear focus on the transformation of representations 
across the thalamocortical bottleneck.

\subsection{Architectural Details}

\subsubsection{RGC Layer}

The model begins with a simulated mosaic of ``ON" and ``OFF" RGCs. Each mosaic of RGCs is arranged in a hexagonal grid with added spatial jitter to mimic the 
quasi-regular sampling observed in the biological retina. Each cell's retina is modeled as a 2D Gaussian filter, where ON-cells are senstive to light and OFF-cells are
sensitive to darkness.

\begin{figure}[h!]
    \centering
    \includegraphics[width=0.9\linewidth]{/Users/sakinkirti/Programming/ucla/dipoppa-lab/02_code/thalamocortical_powerlaw/figures/on_off_RGC.png}
    \caption{Example ON and OFF retinal ganglion cell mosaics.}
    \label{fig:on_off_rgc}
\end{figure}

\subsubsection{LGN Layer}
The LGN layer acts as a relay, roughly maintaining the topology of the RGC input but introducing further stochasticity. In this implementation, the LGN expands the number of 
cells by approximately 2.5x relative to the RGC count. This expansion is achieved by first retaining one copy of each RGC map and further stochastic indexing of 1.5x 
the entire RGC population, creating a redundant but slightly more diverse pool of inputs for the subsequent cortical stage.

\subsubsection{V1 Layer}
The V1 layer is the stage of primary interest, where the massive expansion occurs. Additionally, due to our interest in the connectivity of the V1, we also tested three approaches: 
probabilistic haphazard wiring, unconstrained training, and distance regularized training.

In the probabilistic haphazard wiring approach, we follow the protocol outlined in Ringach 2004, where each V1 cell takes inputs from probabilistically selected LGN neurons
with deterministic synaptic strength. The probability that an LGN-V1 connection exists is modeled as a gaussian function such that as the distance from a V1 receptive field center
increases, the likelihood of an LGN cell forming a connection decreases. Additionally, the synaptic strength between a formalized LGN-V1 connection is given by a gaussian function as 
well.
$$
P=0.85 * \text{exp}\left( -\frac{d}{2\sigma_{\text{conn}}^2} \right), S=1 * \text{exp}\left( -\frac{d}{2\sigma_{\text{syn}}^2} \right)
$$
where $P$ is the probability of a connection, $d$ is the distance, and $S$ is the synaptic strength between a given LGN-V1 cell pair. A major benefit of this paradigm is the 
pre-optimized connectivity matrix that is formed between LGN-V1 pairs, resulting in no need for training. In this probabilistic set up, the V1 layer forms unique receptive 
fields as a result of this haphazard wiring paradigm.

However, it's also reasonable to ask, in a modern deep learning paradigm, if this connectivity (or weights) matrix can be optimized, and if this optimization results in similarly 
formed receptive fields. To this end, we train the same network, retaining the architecture of the RGC and LGN layers (no training), and optimizing the connectivity matrix of the V1
layer in both the unconstrained and distance-regularized settings.

\subsection{Training and benchmarking}

\subsubsection{Optimization procedure}
To train the unconstrained models, we use triplet loss. This objective encourages the model to map ``similar'' images closer together in the representation space than ``dissimilar'' 
images, and is defined as
$$
\mathcal{L}_\text{trip} = \lfloor d(a,p)-d(a,n) + \alpha \rfloor_+
$$
where $a$, $p$, and $n$ denote the anchor, positive, and negative images and $d$ is the distance between two images in representation space. $\alpha$ is a pre-defined margin that allows
the model to make small amounts of errors.

\subsubsection{Benchmarking}
Triplets were generated from natural video datasets, where temporal adjacency defined positive pairs and negative pairs come from pairs of images that are temporally spread apart. 
We used the Adam optimizer with learning rate of 0.0001 and early stopping for 10 epochs based on the validation loss to assess convergence.

To assess representation quality, we use Triplet loss ($\mathcal{L}_{trip}$), violations and Gini index. The violations are defined as the proportion of triplets whose anchor-positive distance is larger
than the anchor-negative distance in representation space, formally given as:
$$
V=
$$

The Gini index is a measure of sparsity, originally developed in economic theories, but reworked here to measure the sparsity of a given V1 representation. We compute the Gini index as:
$$
G=
$$

To measure the overall quality of a particular network architecture, we take some combination of these valuations into account.

\subsection{Implementation and High-Performance Computing}

To handle the scale of these biological simulations, the model was implemented in JAX. This choice allowed for the pre-calculation of connectivity matrices as JAX arrays, 
enabling differentiable forward passes. As V1 dimensions reached the upper bounds of our training paradigm (tens of thousands of neurons), we utilized \textit{multi-GPU sharding} on 
the hoffman2 high performance computing cluster at UCLA. By sharding the V1 connectivity matrix across multiple RTX 2080 Ti GPUs, we bypassed the memory limitations of single devices, 
allowing us to perform full-gradient updates on models that would otherwise be impossible to fit in a single VRAM bank.

\section{Results Summary}

\subsection{Training Convergence}
Models across all V1 scales showed consistent reduction in triplet loss. However, high-expansion models reached lower final losses and exhibited fewer triplet violations on the test set.

\begin{table}[ht!]
\centering
\caption{Model Performance Across Expansion Factors}
\label{tab:results}
\begin{tabular}{@{}lccc@{}}
\toprule
\textbf{V1 Expansion Factor} & \textbf{Final Test Loss} & \textbf{Triplet Violations (\%)} \\ \midrule
1x (Control)        & 0.45            & 18.2\%                   \\
10x                 & 0.32            & 12.5\%                   \\
50x                 & 0.28            & 8.1\%                    \\
300x                & 0.25            & 5.4\%                    \\ \bottomrule
\end{tabular}
\end{table}

\begin{figure}[ht!]
\centering
%\includegraphics[width=0.8\textwidth]{figures/training_curves.png}
\caption{Placeholder: Training curves showing Test Loss vs. Epoch for different V1 dimensions.}
\label{fig:loss}
\end{figure}

\subsection{Representation Sparsity}
Using the \textbf{Gini Coefficient} as a measure of sparsity, we observed that V1 expansion leads to significantly more sparse representations. 

\begin{figure}[ht!]
\centering
%\includegraphics[width=0.8\textwidth]{figures/sparsity_hist.png}
\caption{Placeholder: Histograms of V1 activations and Gini coefficient vs. Expansion Factor.}
\label{fig:sparsity}
\end{figure}

\subsection{V1 receptive fields}
Interestingly, we also identify unique receptive fields that are created through each of the different models. In the haphazard wiring model, we identify unexpected receptive fields
based on the field's prior estimated Gabor filter model of V1 cells.
\begin{figure}[h!]
    \centering
    \includegraphics[width=0.9\linewidth]{/Users/sakinkirti/Programming/ucla/dipoppa-lab/02_code/thalamocortical_powerlaw/figures/haphazard_v1_rf.png}
    \caption{Example receptive fields created by haphazard wiring in V1 layer. ON LGN connections in red, OFF LGN connections in blue.}
    \label{fig:haphazard_v1_rf}
\end{figure}

\section{Discussion and Conclusion}
Our results suggest that the thalamocortical expansion serves as a critical component for metric learning and feature separation. By projecting low-dimensional retinal inputs into a high-dimensional cortical space, the visual system creates a representation where temporal variations are more easily separable. The transition to JAX and multi-GPU sharding has enabled us to test these biological hypotheses at a scale previously reserved for traditional deep learning architectures.

\section{References}
\begin{itemize}
    \item Ringach, D. L. (2002). ``Spatial structure and age-dependent changes in the receptive field properties of V1 neurons.''
    \item Slapik \& Shouval (2026). ``Early Visual Systems: A New Perspective on Expansion.''
    \item JAX Documentation: High-performance gradient transformation.
\end{itemize}

\end{document}
